\honest{My Wild and Reckless Youth}{Eliezer Yudkowsky}{2007}

Parents, it is said, do all the things they tell their children not to do. When I was young and a ``Traditional Rationalist,'' I gave a mysterious answer to a mysterious question. It is the most embarrassing mistake of my life.

I will not tell you what the answer was: ``it would be a long tale and complicated.'' But I can tell you what went wrong. I knew Occam's razor but not the conjunction fallacy, and I thought complicated thoughts without pinning every step down. After my answer, the phenomenon was as mysterious as before.

I was not stupid, I say; the same errors are ``still being made today by respected scientists in respected journals.'' I name none.

I followed the rules. My theory made a bold prediction: neurons would turn out to exploit quantum gravity, ``a la Sir Roger Penrose.'' \nb{The question, which the post does not name, was consciousness. In 2000 the physicist Max Tegmark calculated that such quantum states would decay far too fast to matter, and Hameroff and colleagues disputed the calculation the same year. The idea was argued over by calculation, not left to wait decades for a test.}

But my theory made no retrospective predictions, and ``According to Traditional Science, retrospective predictions don't count.'' A Bayesian would have asked why I believed anything more than ``I don't know.'' \nb{Science has counted old data: the long-known anomaly in Mercury's orbit confirmed general relativity. Explaining why such evidence counts has been a known difficulty for Bayesian theories of confirmation. The post does not put the Bayesian question to the theory it withholds.}

My allegiance to Traditional Rationality, I admit, ``helped me get out of the hole I dug myself into.'' It just was not enough. From my case I conclude: ``You need one whole hell of a lot of rationality before it does anything but lead you into new and interesting mistakes.''

Traditional Rationality, I explain, is learned from the biographies of famous physicists. \nb{The post quotes no text or teacher of it.} It would have let me spend thirty years on my idea. Its followers may agree to disagree, for it lacks the ideal of ``only one correct probability estimate given the evidence.'' \nb{That ideal divides Bayesians as well. Subjective Bayesians permit any coherent starting probabilities, so two of them can disagree on the same evidence.}

The Way of Bayes is also an imprecise art, I concede, but it has ``underlying math, plus experimental evidence from cognitive psychology.'' ``Maybe that will be enough.''
